\section{Introducción}

La optimización de procesos es fundamental en la industria química, debido a que constituye un factor clave para mejorar la eficiencia operativa, reducir costos y mantener una ventaja competitiva en el mercado global. Sin embargo, la búsqueda continua de un mejor rendimiento se enfrenta al desafío de gestionar procesos químicos no lineales y altamente complejos \cite{ref1, ref2}. Los sistemas industriales compuestos por múltiples equipos presentan comportamientos dinámicos complejos caracterizados por fuertes interacciones entre variables, no linealidades y acoplamientos que dificultan su análisis y control \cite{ref3, ref4}. Estas características representan un desafío significativo para los métodos de control convencionales, particularmente aquellos basados en lazos de control independientes, como los controladores PID (Proporcional–Integral–Derivativo), los cuales pueden presentar limitaciones en el manejo de dichas interacciones, así como sensibilidad a perturbaciones y efectos de retardo. Como consecuencia, pueden generarse respuestas lentas, sobreoscilaciones y dificultades para mantener la estabilidad y el desempeño óptimo del proceso. Estas características pueden generar respuestas lentas, sobreoscilaciones y dificultades para mantener la estabilidad y el desempeño óptimo del proceso.Con el fin de superar estas limitaciones, el Control Predictivo Basado en Modelo (Model Predictive Control, MPC) se ha consolidado como una de las estrategias de control avanzado más utilizadas en procesos industriales. Este método emplea un modelo dinámico del proceso para predecir su comportamiento futuro dentro de un horizonte de tiempo determinado y resolver, en cada instante de operación, un problema de optimización que permite determinar las acciones de control más adecuadas \cite{ref5}. Gracias a su capacidad para manejar sistemas multivariables, considerar restricciones operativas y anticipar el efecto de perturbaciones, el MPC ha sido ampliamente adoptado en diversos sectores industriales. No obstante, incluso las estrategias de control avanzadas, como el MPC, presentan dificultades cuando la dinámica del proceso no se conoce con precisión o está sujeta a variaciones significativas. La complejidad numérica, caracterizada por fuertes interacciones entre variables del proceso y dinámicas no modeladas, junto con la necesidad de resolver continuamente problemas de optimización, puede traducirse en respuestas insuficientes o en la falta de convergencia del controlador. Como respuesta a estas limitaciones, en los últimos años ha surgido el uso de modelos subrogados, también conocidos como metamodelos, los cuales buscan aproximar el comportamiento de modelos rigurosos mediante representaciones matemáticas más simples construidas a partir de datos \cite{ref6, ref7}. Estos modelos permiten reducir significativamente la carga computacional asociada a la simulación y optimización de procesos complejos, manteniendo al mismo tiempo un nivel adecuado de precisión en la predicción del comportamiento del sistema \cite{ref8, ref9}. El presente trabajo se enmarca en la continuidad de una línea de investigación en la cual se desarrolló previamente la subrogación de control para un proceso químico complejo con el objetivo de implementar estrategias de control predictivo orientadas al rechazo de perturbaciones en las variables del sistema \cite{ref10}. En dicho estudio se utilizó un modelo subrogado basado en representación en espacio de estados, el cual será empleado en este trabajo como referencia de comparación frente a otros métodos de identificación y modelado, específicamente los modelos ARX, ARMAX, de primer orden con tiempo muerto (FOPDT) y modelos basados en función de transferencia.
